\begin{center}
\begin{tikzpicture}[
  >=Latex,
  every node/.style={font=\small},
  box/.style={draw,rounded corners=2pt,align=center,inner sep=6pt,
             text width=3.15cm,minimum height=1.02cm},
  flow/.style={->,thick}
]
\node[font=\normalsize\bfseries] at (6.0,2.65)
  {超越似然：概览};

\node[box] (gan) at (1.7,0.85) {
  \textbf{GAN 比较}\\[2pt]
  生成器样本\\
  $\leftrightarrow$ 判别器\\
  反馈
};
\node[box] (wgan) at (6.0,0.85) {
  \textbf{WGAN 几何}\\[2pt]
  传输距离\\
  取代\\
  分类
};
\node[box] (ebm) at (10.3,0.85) {
  \textbf{EBM 景观}\\[2pt]
  能量赋予\\
  相对\\
  合理性
};
\draw[flow] (gan) -- (wgan);
\draw[flow] (wgan) -- (ebm);

\node[box] (ratio) at (1.7,-2.15) {
  \textbf{最优}\\[-1pt]
  \textbf{判别器}\\[2pt]
  局部密度比\\
  信息
};
\node[box] (energy) at (6.0,-2.15) {
  \textbf{未归一化密度}\\[2pt]
  $p_\theta(x)\propto e^{-E_\theta(x)}$\\
  全局\\
  配分函数
};
\node[box] (score) at (10.3,-2.15) {
  \textbf{回到得分}\\[2pt]
  $\nabla_x\log p_\theta(x)$\\
  $=-\nabla_xE_\theta(x)$
};
\draw[flow] (gan) -- (ratio);
\draw[flow] (ebm) -- (energy);
\draw[flow] (energy) -- (score);
\end{tikzpicture}
\end{center}